%% ============================================================================
%%  The Prefetch Paradox That Wasn't: Four Measurement Artifacts in a
%%  CRIU Lazy-Restore Study.  IEEEtran conference format (venue TBD).
%%  Build: pdflatex pitfall && bibtex pitfall && pdflatex pitfall && pdflatex pitfall
%% ============================================================================
\documentclass[conference]{IEEEtran}

\usepackage[T1]{fontenc}
\usepackage[utf8]{inputenc}
\usepackage{microtype}
\usepackage{amsmath,amssymb}
\usepackage{booktabs}
\usepackage{multirow}
\usepackage{xcolor}
\usepackage{pgfplots}
\pgfplotsset{compat=1.18}
\usepackage{cite}
\usepackage[colorlinks=true,linkcolor=black,citecolor=black,urlcolor=blue!60!black]{hyperref}
\usepackage{url}
\usepackage{xspace}

% Categorical slots 1-4 (validated: CVD-safe adjacent pairs; markers + table
% carry identity for the two sub-3:1 hues).
\definecolor{cA}{HTML}{2A78D6}
\definecolor{cB}{HTML}{EB6834}
\definecolor{cC}{HTML}{1BAF7A}
\definecolor{cD}{HTML}{EDA100}

\newcommand{\uffd}{\texttt{userfaultfd}\xspace}
\newcommand{\nd}{\texttt{TCP\_NODELAY}\xspace}
\newcommand{\us}{\,$\mu$s\xspace}

\title{The Prefetch Paradox That Wasn't:\\
Four Measurement Artifacts in a CRIU Lazy-Restore Study}

\author{\IEEEauthorblockN{Utkarsh Maurya}
\IEEEauthorblockA{Independent Researcher\\
\href{mailto:utkarsh@kernex.sbs}{utkarsh@kernex.sbs}}}

\begin{document}
\maketitle

\begin{abstract}
We had a near-submission paper reporting that fixed sequential prefetch
\emph{slows} post-restore time-to-first-request (TTFR) under CRIU lazy restore
by 88\,\% at loopback, that an RTT crossover near 100--150\us separates a
harmful regime from a helpful one, and that an adaptive controller recovers the
loss. Every one of these findings is an artifact. Nagle's algorithm on the page
server, interacting with delayed ACKs, stalls each pipelined prefetch batch by
${\approx}41$\,ms. The stall backs up the prefetch queue, and a linear
duplicate scan on the fault path turns that backlog into CPU time charged to
every fault. A teardown race silently recorded lost statistics as zeros,
fabricating an 85\,\% fault reduction. A probe race inflated full-restore TTFR
of a small workload 60$\times$, inventing a 24$\times$ lazy-restore win that is
in fact a 2.4$\times$ loss. In a single session we reproduce the original
numbers within 10\,\% with Nagle on and invert them with Nagle off: prefetch
cuts TTFR by 48\,\% at loopback and 87\,\% at 1\,ms RTT, with no crossover
anywhere from 0 to 2\,ms. We describe how each artifact hid, the signature
that would have exposed it, and a short checklist for remote-paging
measurements.
\end{abstract}

\begin{IEEEkeywords}
measurement pitfalls, TCP Nagle algorithm, delayed ACK, lazy restore, CRIU,
userfaultfd, prefetching, reproducibility
\end{IEEEkeywords}

% ─────────────────────────────────────────────────────────────────────────────
\section{Introduction}

A counter-intuitive result is the best kind to publish and the worst kind to
trust. Our earlier study of CRIU lazy restore~\cite{criu,criu-lazy} found
one: speculative prefetch of neighboring pages, which cut fault counts by
85\,\%, nonetheless made a PyTorch process 88\,\% slower to answer its first
request. We explained it as contention between prefetch traffic and fault
resolution, swept emulated RTT to locate a crossover where prefetch turns
helpful, confirmed the helpful side on a two-machine LAN, and designed an
adaptive controller that detects the ``congested'' regime and backs off. The
results were statistically sharp ($n\!=\!20$, Welch $|t|>45$) and internally
consistent.

They were also wrong. While building a successor system we found its bandwidth
estimator reporting ${\sim}700$ pages/s on loopback, two orders of magnitude
below what the channel should carry. That number led to Nagle's
algorithm~\cite{rfc896}; pulling that thread unraveled the paper. This paper
documents what we found, because each artifact is generic to remote-paging
prototypes and none announced itself:

\begin{enumerate}
\item \textbf{Nagle $\times$ delayed ACK} on a request/response protocol with
pipelined sub-MSS responses stalls every prefetch batch by one delayed-ACK
timeout (\S\ref{sec:nagle}).
\item \textbf{A fault-path coupling}: the stalled queue grows, and an
$O(\text{depth})$ duplicate scan executed per fault converts it into a fixed
CPU tax, the ``contention'' we had attributed to the network
(\S\ref{sec:coupling}).
\item \textbf{Silent zeros}: when the stalled worker outlived the harness's
kill deadline, its statistics were lost and parsed as zero, fabricating the
85\,\% fault reduction that made the result paradoxical (\S\ref{sec:zeros}).
\item \textbf{A probe race}: a passive TTFR probe deleted its signal file after
restore began, so a fast restore was charged a full one-second sleep tick
(\S\ref{sec:probe}).
\end{enumerate}

\noindent We confirm each with a targeted A/B on the same machine, kernel and
build in one session (\S\ref{sec:corrected}). With all four fixed, prefetch
helps at every RTT we measured, and the controller's ``recovery'' turns out to
have been a response to symptoms of the transport bug (\S\ref{sec:controller}).
\S\ref{sec:lessons} distills a checklist.

% ─────────────────────────────────────────────────────────────────────────────
\section{System and Original Claims}
\label{sec:system}

\textbf{Lazy restore.}\, CRIU's lazy restore resumes a checkpointed process
with its memory unpopulated; a handler registered on \uffd~\cite{userfaultfd}
receives each page fault, fetches the page from a page server and installs it
with \texttt{UFFDIO\_COPY}. This is post-copy migration~\cite{hines2009postcopy}
at process granularity, and the same TTFR concern drives snapshot work such as
REAP~\cite{ustiugov2021reap}.

\textbf{Our prototype.}\, A C handler resolves faults \emph{serially} on one
TCP connection: an 8-byte address request, a 4096-byte page response. On
each fault it also enqueues up to 16 sequential (or 8 strided) neighbors for a
background worker, which pulls batches of up to 32 addresses from the queue,
sends all 32 requests back-to-back on a second connection, then reads the 32
responses. Before enqueuing, the fault path scans the queue for duplicates. A
Python page server serves each connection on its own thread.

\textbf{Measurement.}\, TTFR is the time from invoking \texttt{criu restore} to
the workload's first observable response. Workloads: \emph{test\_loop}, a C
loop that rewrites a counter file once per \texttt{sleep(1)}; Valkey (Redis)
with 10k keys; and a ResNet-18 PyTorch process (${\approx}15.5$k pages touched
before first inference). Modes: \texttt{full} restore, demand-only
\texttt{lazy}, fixed \texttt{prefetch}, and \texttt{adaptive} (prefetch with a
controller that halves or disables prefetch on high duplicate rate or queue
depth). RTT is injected with \texttt{netem}~\cite{hemminger2005netem} on
loopback.

\textbf{Original claims} (all PyTorch unless noted, $n\!=\!20$):
\begin{itemize}
\item[C0] Prefetch \emph{hurts} below an RTT crossover of 100--150\us and helps
above it: $+88\,\%$ at loopback ($650\to1227$\,ms), $-37\,\%$ at 1\,ms.
\item[C1] Lazy restore cuts test\_loop TTFR $24\times$ (full $1019$\,ms vs
lazy $42$\,ms).
\item[C2] Prefetch hurts because prefetch traffic contends with fault
resolution, even while reducing faults by 85\,\% ($15{,}515\to2{,}322$).
\item[C3] The adaptive controller recovers the regression ($655$\,ms, within
noise of lazy).
\end{itemize}

% ─────────────────────────────────────────────────────────────────────────────
\section{Artifact 1: Nagle Meets Delayed ACK}
\label{sec:nagle}

\textbf{Mechanism.}\, Nagle's algorithm~\cite{rfc896} holds a sub-MSS segment
while earlier data on the connection is unacknowledged. Delayed
ACK~\cite{rfc1122} lets a receiver with nothing to send wait before
acknowledging; Linux waits at least 40\,ms. Their interaction is a known
hazard for request/response traffic~\cite{minshall2000pitfalls,
mogul2001rethinking}, and our prefetch protocol walks straight into it. The
worker sends 32 requests; the server answers the first with a 4096-byte
segment, which is below the loopback MSS (${\approx}64$\,KB) and also leaves a
sub-MSS tail on an Ethernet MSS. The second response is held by Nagle until
that segment is acknowledged. The client has already sent every request and
has nothing to piggyback an ACK on, so it delays: ${\approx}41$\,ms per batch.
The \emph{fault} connection is immune, because a serial exchange never has
unacknowledged data when the next response is written. Prefetch was the only
path that pipelined, so it was the only path that stalled.

\begin{table}[t]
\centering
\caption{Standalone repro of the page-server protocol on loopback (no netem):
median batch time and throughput, 200 batches. Only the \emph{server}'s
\nd matters.}
\label{tab:repro}
\small
\begin{tabular}{@{}rrrrr@{}}
\toprule
& \multicolumn{4}{c}{\nd (server, client)} \\
\cmidrule(l){2-5}
batch & off, off & on, off & off, on & on, on \\
\midrule
1  & 18\us      & 15\us     & 17\us      & 16\us \\
8  & 41{,}009\us & 77\us    & 41{,}009\us & 62\us \\
32 & 40{,}893\us & 415\us   & 41{,}096\us & 288\us \\
\midrule
32 (pages/s) & 1{,}092 & 71{,}849 & 850 & 105{,}267 \\
\bottomrule
\end{tabular}
\end{table}

\textbf{Evidence.}\, Table~\ref{tab:repro} reproduces the stall outside the
system with a short script mimicking the page server and the worker
(\texttt{eval/nagle\_repro.py}, no root). A single-page exchange is
unaffected. Any batch of $\ge2$ costs one delayed-ACK timeout, and setting
\nd on the server alone removes it. The worker therefore moved roughly
$32/41\,\text{ms}\approx780$ pages/s, which is the figure our successor's
estimator had reported.

\textbf{Signature.}\, A latency quantum near 40\,ms, independent of
payload size, that appears only with pipelining. We never looked, because
the paper's metric was end-to-end TTFR, in which the quantum was smeared
across thousands of batches.

% ─────────────────────────────────────────────────────────────────────────────
\section{Artifact 2: How a Stalled Side Channel Slows the Main Path}
\label{sec:coupling}

Nagle explains why prefetch was useless at loopback. It does not explain
why prefetch was \emph{harmful}, since the serial fault path never stalls.
Yet TTFR rose by ${\approx}600$\,ms. The coupling is in the handler, not the
network. On every fault, \texttt{async\_prefetch\_enqueue} takes the queue
lock and, for each candidate, scans the whole queue to drop duplicates. The
worker drains ${\approx}780$ pages/s while each fault offers up to 16
candidates, so the queue fills to its 8{,}192-entry cap within the first
second and stays there, dropping ${\approx}124$k requests per run. The build
used \texttt{-O0}.

\textbf{Causal test.}\, We instrumented the handler to time enqueue on the
fault path and record queue depth at each enqueue, and added a switch that
skips the scan. Table~\ref{tab:ablate} is a 2$\times$2 of \nd $\times$ scan
at loopback. With Nagle on, removing the scan alone takes prefetch from
$1207$ to $608$\,ms, indistinguishable from lazy ($601$\,ms). The enqueue time
it removes, 577\,ms or 37\us per fault, accounts for the whole regression.
Fault counts do not move (15.5k in every Nagle-on cell), so prefetch was
neither helping nor competing for the wire; it was burning the fault path's
CPU. With \nd the queue averages 235 entries, the scan costs 24\,ms in total,
and keeping it is slightly \emph{better} ($299$ vs $323$\,ms): without
deduplication the worker wastes bandwidth on repeated pages and faults rise
from 5.0k to 6.1k.

\begin{table}[t]
\centering
\caption{Causal test, PyTorch fixed prefetch at loopback, $n\!=\!10$:
TTFR (mean $\pm$ 95\,\% CI), page faults, total enqueue time on the fault
path, and mean queue depth at enqueue (cap 8{,}192).}
\label{tab:ablate}
\small
\setlength{\tabcolsep}{4pt}
\begin{tabular}{@{}llrrrr@{}}
\toprule
\nd & dup.\ scan & TTFR (ms) & faults & enq.\ (ms) & depth \\
\midrule
off & linear  & $1207\pm33$ & 15.5k & 577 & 7{,}657 \\
off & skipped & $608\pm15$  & 15.5k & 1   & 8{,}048 \\
on  & linear  & $299\pm7$   & 5.0k  & 24  & 235 \\
on  & skipped & $323\pm10$  & 6.1k  & 1   & 1{,}203 \\
\midrule
off & lazy    & $601\pm9$   & 15.5k & --  & -- \\
on  & lazy    & $584\pm7$   & 15.5k & --  & -- \\
\bottomrule
\end{tabular}
\end{table}

\begin{table}[t]
\centering
\caption{Microbenchmark of the duplicate scan (16 candidates, all misses)
against queue depth, scaled by 15.5k faults: the debug build costs
${\approx}4\times$.}
\label{tab:scan}
\small
\begin{tabular}{@{}rrrrr@{}}
\toprule
& \multicolumn{2}{c}{\texttt{-O0} (as built)} & \multicolumn{2}{c}{\texttt{-O2}} \\
\cmidrule(lr){2-3}\cmidrule(l){4-5}
depth & \us/fault & ms total & \us/fault & ms total \\
\midrule
512   & 7.7   & 119   & 1.9  & 30  \\
1{,}500 & 22.2 & 344  & 5.5  & 85  \\
3{,}000 & 44.9 & 696  & 11.2 & 173 \\
6{,}000 & 90.6 & 1{,}404 & 22.2 & 343 \\
\bottomrule
\end{tabular}
\end{table}

Table~\ref{tab:scan} isolates the scan in a microbenchmark. Its worst case
over-predicts the measured cost, because at a full queue each enqueue scans
for the first candidate and then drops the rest, but it shows the build flag
multiplying the tax by about four. Nagle supplied the backlog, the scan turned
the backlog into fault-path latency, and the debug build amplified it.

\textbf{The crossover, explained.}\, With Nagle on, the stalled worker still
lands pages ahead of demand once each fault is slow enough. Its fault savings
grow with RTT: PyTorch faults under prefetch fall from 15.4k at loopback to
14.4k at 150\us, 13.1k at 200\us and 9.6k at 1\,ms. Each saved fault is worth
one lazy per-fault cost (lazy TTFR$/15.5$k: 205\us at 150\us RTT, 253\us at
200\us). Savings therefore reach the ${\approx}550$\,ms scan tax between 150
and 200\us. Measured Nagle-on prefetch breaks even at 150\us (ratio 1.00) and
wins by 12\,\% at 200\us. The ``crossover'' we reported as a property of the
network was the point where a fixed self-inflicted tax met RTT-proportional
savings.

% ─────────────────────────────────────────────────────────────────────────────
\section{Artifact 3: Silent Zeros}
\label{sec:zeros}

The 85\,\% fault reduction made C2 paradoxical, and the paradox made it
interesting. It was never real. The harness sends \texttt{SIGTERM} to the
handler after measuring, waits 3\,s, then \texttt{SIGKILL}s it and parses its
log for a statistics line. On stop, the worker drained its queue first. With
Nagle on, a full 8{,}192-entry queue at ${\approx}780$ pages/s needs about
10\,s,
the handler died before printing, and the parser's defaults (all zero) went
into the CSV. In the original data, 17 of 20 prefetch rows read zero faults;
the other three read 15{,}440--15{,}498, the same as lazy. Our rerun
reproduces the pattern: 16 of 20 zeros, and 15{,}452--15{,}650 in the four
complete rows. The mean of a mixture of zeros and real values produced
``$2{,}322$''.

Prefetch with Nagle on reduced faults by essentially \emph{nothing}. The
tell was in plain sight: a per-mode mean whose rows were bimodal between 0 and
the lazy value. We fixed the worker to exit on stop without draining, and the
parser now writes $-1$ (missing) rather than 0. With the fix, every
Nagle-on prefetch row in Table~\ref{tab:ablate} reports its faults:
15.4--16.0k.

% ─────────────────────────────────────────────────────────────────────────────
\section{Artifact 4: A Probe Race}
\label{sec:probe}

test\_loop's TTFR probe is passive: it deletes the counter file and waits for
the restored process to rewrite it. The deletion happened \emph{after} restore
began. A fast restore (full restore of a 1\,MB process) often rewrote the file
before the probe deleted it, so the probe waited for the next write, one
\texttt{sleep(1)} later. Full restore was charged $1019\pm1$\,ms. Lazy restore
was slow enough to lose the race and was charged its true ${\approx}42$\,ms.
Moving the deletion to the window between checkpoint and restore gives full
$17$\,ms and lazy $41$--$42$\,ms. Lazy restore is $2.4\times$ \emph{slower},
not $24\times$ faster.

\textbf{Signature.}\, A result that equals a known period (1000\,ms here,
${\approx}$ one tick of \texttt{sleep(1)}) plus a small constant, with a
suspiciously tight CI ($\pm1$\,ms over 20 runs). Full restore of 266 pages at
1\,s also implies 1\,MB/s, three orders of magnitude below a memory copy.

% ─────────────────────────────────────────────────────────────────────────────
\section{Corrected Results}
\label{sec:corrected}

\textbf{Setup.}\, AMD Ryzen 7 7735HS (8 cores, 16 threads), 14\,GB RAM, Linux
7.2.8, CRIU 4.2.1, Python 3.14, PyTorch 2.15 (CPU). The original paper used
the same hardware on Linux 6.18.7 with PyTorch 2.12. Each configuration runs
with \nd off (the original transport) and on, arms interleaved per workload
and per RTT point, in one session, with the probe race fixed in both arms.
Matrix: $n\!=\!20$ at loopback. Sweep: PyTorch, $n\!=\!10$, RTT 0--2000\us.
Values are mean $\pm$ 95\,\% CI.

\begin{table}[t]
\centering
\caption{Loopback TTFR (ms), $n\!=\!20$. ``Paper'' is the original dataset
(Nagle on, probe race present); the Nagle-on rerun reproduces it.}
\label{tab:matrix}
\small
\setlength{\tabcolsep}{4pt}
\begin{tabular}{@{}llrrr@{}}
\toprule
workload & mode & paper & Nagle on & \nd \\
\midrule
\multirow{4}{*}{PyTorch}
 & full     & $191\pm7$  & $202\pm5$  & $208\pm2$ \\
 & lazy     & $650\pm10$ & $618\pm8$  & $618\pm16$ \\
 & prefetch & $1227\pm23$ & $1146\pm15$ & $\mathbf{311\pm11}$ \\
 & adaptive & $655\pm9$  & $594\pm7$  & $352\pm55^{*}$ \\
\midrule
\multirow{4}{*}{test\_loop}
 & full     & $1019\pm1^{\dagger}$ & $17\pm0$ & $18\pm0$ \\
 & lazy     & $42\pm1$   & $41\pm1$ & $42\pm2$ \\
 & prefetch & $43\pm3$   & $41\pm1$ & $27\pm1$ \\
 & adaptive & $41\pm1$   & $40\pm2$ & $29\pm1$ \\
\midrule
\multirow{4}{*}{Redis}
 & full     & $32\pm2$ & $31\pm1$ & $31\pm0$ \\
 & lazy     & $37\pm2$ & $35\pm1$ & $35\pm2$ \\
 & prefetch & $39\pm2$ & $37\pm2$ & $37\pm1$ \\
 & adaptive & $38\pm2$ & $34\pm1$ & $34\pm1$ \\
\bottomrule
\multicolumn{5}{@{}p{0.95\columnwidth}@{}}{\footnotesize $^{*}$One
iteration at 878\,ms; median 325\,ms. $^{\dagger}$Probe race
(\S\ref{sec:probe}).}
\end{tabular}
\end{table}

\textbf{Attribution.}\, With Nagle on, the rerun reproduces the paper on a
different kernel and PyTorch version: prefetch $+85\,\%$ vs lazy (paper:
$+89\,\%$), adaptive within 4\,\% of lazy. The only remaining
difference between the Nagle-on and \nd columns is one socket option, and
it flips PyTorch prefetch from $+85\,\%$ to $-50\,\%$ ($p\ll10^{-6}$). Redis
is unaffected either way, since its first request touches few pages.
test\_loop's prefetch gain ($-35\,\%$) appears only with \nd.

\begin{figure}[t]
\centering
\begin{tikzpicture}
\begin{axis}[
  width=\columnwidth, height=5.6cm,
  symbolic x coords={0,30,60,100,150,200,300,500,1000,2000},
  xtick={0,30,60,100,150,200,300,500,1000,2000}, xticklabel style={font=\scriptsize},
  xlabel={injected RTT (\textmu s, non-uniform spacing)},
  ymode=log, ylabel={TTFR (ms)}, ymin=200, ymax=30000,
  yticklabel style={font=\scriptsize}, label style={font=\small},
  ymajorgrids, major grid style={black!12},
  legend style={font=\scriptsize, at={(0.02,0.98)}, anchor=north west,
                draw=none, fill=white, fill opacity=0.85, text opacity=1},
  legend cell align=left,
  every axis plot/.append style={line width=1pt, mark size=2pt},
  unbounded coords=jump,
]
\addplot[color=cD, mark=square*] table[x=rtt, y=lazy]{sweep.dat};
\addlegendentry{lazy (demand only)}
\addplot[color=cB, mark=triangle*] table[x=rtt, y=pf0]{sweep.dat};
\addlegendentry{prefetch, Nagle on}
\addplot[color=cA, mark=*] table[x=rtt, y=pf1]{sweep.dat};
\addlegendentry{prefetch, \nd}
\addplot[color=cC, mark=diamond*] table[x=rtt, y=ad1]{sweep.dat};
\addlegendentry{adaptive, \nd}
\end{axis}
\end{tikzpicture}
\caption{PyTorch TTFR vs injected RTT ($n\!=\!10$ per point). With Nagle on,
prefetch is slower than lazy up to 150\us; with \nd it is 2--8$\times$
faster at every RTT. Lazy exceeds the 30\,s probe timeout at 2\,ms.}
\label{fig:sweep}
\end{figure}

\textbf{No crossover.}\, Figure~\ref{fig:sweep} and Table~\ref{tab:sweep}
show the sweep. The Nagle-on arm reproduces the original curve, breaking even
at 150\us. With \nd, prefetch beats lazy at every RTT, and the margin widens
with RTT, from $0.52\times$ lazy's TTFR at loopback to $0.13\times$ at 1\,ms.
Faults under prefetch fall from 5.0k at loopback to a floor of 2.0k from
200\us on, against 15.5k for lazy. The original paper's one surviving claim,
that prefetch helps at high RTT, holds and is several times stronger
($-87\,\%$ at 1\,ms, versus the reported $-37\,\%$).

\begin{table}[t]
\centering
\caption{PyTorch sweep, TTFR relative to each arm's own lazy at the same RTT ($n\!=\!10$); lazy row from the \nd arm.}
\label{tab:sweep}
\scriptsize
\setlength{\tabcolsep}{2.2pt}
\begin{tabular}{@{}lrrrrrrrrr@{}}
\toprule
RTT\,(\textmu s) & 0 & 30 & 60 & 100 & 150 & 200 & 300 & 500 & 1000 \\
\midrule
lazy (ms) & 638 & 1194 & 1606 & 2270 & 3171 & 3918 & 5587 & 8788 & 16672 \\
\midrule
pf, Nagle on & 1.84 & 1.50 & 1.44 & 1.10 & 1.00 & 0.88 & 0.82 & 0.74 & 0.62 \\
pf, \nd      & 0.52 & 0.34 & 0.26 & 0.20 & 0.17 & 0.15 & 0.14 & 0.14 & 0.13 \\
ad, \nd      & 0.51 & 0.34 & 0.25 & 0.24 & 0.16 & 0.15 & 0.13 & 0.11 & 0.10 \\
\bottomrule
\end{tabular}
\end{table}

% ─────────────────────────────────────────────────────────────────────────────
\section{The Controller Learned the Bug}
\label{sec:controller}

C3's controller disables prefetch when the duplicate rate exceeds 80\,\% or
the queue is deep. With Nagle on, its first 128-fault window saw 48 pages
prefetched, a 92\,\% duplicate rate and a depth of 1{,}381, and it switched
prefetch off. That is the right response to a stalled channel, and it is why
adaptive matched lazy. But the condition it detected was the bug, not a
property of the workload or the network. Signals tuned on an artifact can
look validated, because the artifact is exactly what they react to.

With \nd, the same controller shrinks its window instead of disabling it,
and lands slightly \emph{ahead} of fixed prefetch at high RTT ($-22\,\%$ at
1\,ms, $-24\,\%$ at 2\,ms). We report this without claiming a mechanism; it
is the starting point for a corrected study, not a result of this one. Our
successor system's ``bandwidth-aware'' regime gate had the same flaw: it was
estimating the Nagle-limited ${\sim}700$ pages/s and switching regimes on it.

% ─────────────────────────────────────────────────────────────────────────────
\section{Lessons}
\label{sec:lessons}

\textbf{Treat a paradox as a bug until the mechanism is measured, not
inferred.} We had a plausible story (contention), consistent aggregate data
and large $t$ statistics. None of it tested the story. A per-component time
breakdown (fault-path time per fault with and without prefetch) would have
shown the cost sitting in the handler's CPU, not on the wire.

\textbf{Set \nd on any request/response socket that pipelines, and test with
a two-request microbenchmark.} The hazard is decades old~\cite{minshall2000pitfalls,
mogul2001rethinking}, but it bites only when a code path pipelines, so a
serial smoke test passes. Look for 40\,ms quanta.

\textbf{Loopback is not a neutral baseline.} Loopback's ${\approx}64$\,KB MSS
makes every page a ``small'' segment for Nagle, so it maximized the artifact
in exactly the configuration we called best-case.

\textbf{Never let missing data default to zero.} Emit a sentinel, count
missing rows per cell, and look at per-row distributions before means. A
bimodal cell (zeros mixed with real values) is a pipeline bug.

\textbf{Distrust results that equal a constant of the setup.} 1019\,ms on a
workload that sleeps for one second; 41\,ms batches on a kernel whose
minimum delayed ACK is 40\,ms; ${\sim}700$ pages/s that equals $32/41$\,ms.
Each artifact had a number that matched a constant of the setup.

\textbf{Reproduce the old result before claiming the new one.} Running both
arms interleaved in one session, and recovering the original numbers in the
``bug on'' arm, attributes the change to the bug rather than to a new
kernel, library version or machine state.

These echo general warnings about measurement bias~\cite{mytkowicz2009wrong}
and benchmarking practice~\cite{hoefler2015scientific,vanderkouwe2019sok};
our contribution is a concrete, end-to-end case in remote paging where four
artifacts compounded into a coherent, wrong story.

% ─────────────────────────────────────────────────────────────────────────────
\section{Limitations}

All corrected measurements are single-host with netem-injected RTT. The
original two-machine LAN result ($-72\,\%$ at 311\us) ran with Nagle on and
has not been redone; we expect a larger gain, but have not measured it. One
machine, one kernel and three workloads; the PyTorch adaptive cells carry
single-iteration outliers. The duplicate-scan tax depends on queue depth and
build flags, so its magnitude is specific to our prototype, but the pattern
(a stalled side channel taxing the critical path through shared state) is not.

\section{Artifact Availability}

Code, raw CSVs and per-run handler logs:
\url{https://github.com/kernex-sbs/DistriProc} (branch
\texttt{nodelay-remeasure}). \texttt{eval/remeasure\_nodelay.sh} runs both
arms; \texttt{eval/analyze\_nodelay.py} regenerates Tables~\ref{tab:matrix}
and~\ref{tab:sweep}; \texttt{eval/nagle\_repro.py} reproduces
Table~\ref{tab:repro} without root; \texttt{eval/ablate\_dedup.sh} runs
Table~\ref{tab:ablate}. The handler and page server read
\texttt{DISTRIPROC\_TCP\_NODELAY=0} to restore the original transport.

\bibliographystyle{IEEEtran}
\bibliography{references}

\end{document}
